\documentclass[../main.tex]{subfiles}
\graphicspath{{\subfix{../images/}}}
\begin{document}
\newpage
\section*{Conclusion Générale}
\addcontentsline{toc}{section}{Conclusion Générale}

Ce projet, réalisé dans le cadre du développement de Watto, une plateforme conçue en pensant au marché tunisien, avait pour objectif de révolutionner la manière dont les individus s'engagent dans les cadeaux et les transactions financières. En fournissant une boîte à outils polyvalente pour les cadeaux collectifs, une gestion flexible des fonds et l'intégration de partenariats avec des vendeurs locaux, Watto visait à créer une plateforme à la fois inclusive et innovante.

\vspace{10pt}
\newlineLe projet a suivi une approche structurée, en utilisant CircleCI pour rationaliser le processus d'intégration continue/déploiement continu (CI/CD), en s'intégrant avec Vercel et Jira, et en exploitant des technologies telles que Next.js et Nest.js pour offrir une solution robuste.

\vspace{10pt}
\newlineÀ mesure que nous avançons, les perspectives futures de Watto sont prometteuses. Notre attention reste focalisée sur l'amélioration de l'expérience utilisateur grâce à des recommandations personnalisées et l'extension de la plateforme de défis de codage pour prendre en charge une gamme plus large de langages de programmation et de scénarios. Avec une base solide et un engagement envers l'amélioration continue, Watto est prêt à enrichir davantage l'expérience des cadeaux pour les individus tout en contribuant à la croissance de l'écosystème technologique tunisien.
\end{document}